%% ma: L11-B14-0028
%% lop: 11
%% chuong: 4
%% bai: 14
%% dang: 11.14.4
%% loai: TLN
%% muc: VD
%% dap_an: "0,5"
%% nguon: "tai-lieu-hoc-tap-mon-toan-lop-11-theo-chuong-trinh-giao-duc-pho-thong-2018.pdf (Lớp Toán Thầy Duy, Phạm Lê Duy, THPT Chu Văn An), Chương 4 Bài 5, Phần 3, Câu 4 (tr.535)"
%% kien_thuc: "Lớp 11 Bài 14 (phép chiếu song song)"
%% trang_thai: chinh-thuc
%% ghi_chu: "Đã duyệt 03/11/2026."
\begin{ex}
Cho hình lăng trụ $ABC.A^{\prime}B^{\prime}C^{\prime}$, gọi $I$ là trung điểm $AB$. Gọi $S,S^{\prime}$ lần lượt là diện tích của tam giác $ABC$ và hình chiếu của tam giác $AA^{\prime}C^{\prime}$ lên mặt phẳng $(A^{\prime}B^{\prime}C^{\prime})$ theo phương $IB^{\prime}$. Khi đó tỉ số $\dfrac{S^{\prime}}S$ bằng
\shortans{0,5}
\loigiai{\sloppy Theo phương $IB^{\prime}$: $A\mapsto$ trung điểm $A_1$ của $A^{\prime}B^{\prime}$ (vì $AIB^{\prime}A_1$ là hình bình hành: $\overrightarrow{AA_1}=\overrightarrow{IB^{\prime}}$), $A^{\prime}\mapsto A^{\prime}$, $C^{\prime}\mapsto C^{\prime}$. Hình chiếu là $\triangle A_1A^{\prime}C^{\prime}$ có diện tích $\dfrac12S_{A^{\prime}B^{\prime}C^{\prime}}=\dfrac12S$. Vậy $\dfrac{S^{\prime}}S=0{,}5$.}
\end{ex}
